\section{Introduction}\label{sec:introduction}

The transcendence degree of Siegel's $E$-functions \cite{Siegel1929}
is preserved under specialization at nonzero algebraic ordinary points \cite{Shidlovskii1989}.
Beukers lifts relations \cite[Theorems~1.1--1.3]{Beukers2006}
using Andr\'e's $E$-operators \cite{Andre2000I,Andre2000II}.

Parameter coincidences, Wronskians and shared logarithmic companions
produce mixed relations. We determine their ideal, retaining normalized
additive constants. Nonclosed lists are treated by closure, specialization and elimination below.

\subsection{A first example}
Take $c\ne0$ and the linked blocks $\operatorname{Bess}(c,0)$ and
$\operatorname{Kum}(2c,1/2)$. They form one linkage class and leave three independent coordinates
over $T=\C(z)(e^{cz})$ and four over $\C(z)$.
Supplementary Section~\ref{sup:first-example} gives the maps, determinant, series and connections.

\subsection{The main results}\label{subsec:normalization}
The following data conventions define all terms in Theorem~\ref{thm:main}.
Work over $k=\C(z)$ with $D=d/dz$ and a finite admissible coordinate list
$\mathcal F$, including repeated input slots. Fix a simply connected
$\Omega_{\rm an}\subset\C^\times$ avoiding the finite poles of the data,
a point $z_*\in\Omega_{\rm an}$ and a branch $L=\log z$ there.
All functions are meromorphic on this domain, with differential constant
field $\C$. Set $\phi_{\mu,p}=e^{\mu(z-p)}/(z-p)$,
$\Phi'_{\mu,p}=\phi_{\mu,p}$, $\Phi_{\mu,p}(z_*)=0$ for
$(\mu,p)\ne(0,0)$, and $\Phi_{0,0}=L$.
A prime on a coordinate denotes $D$; on a function of its own argument
it denotes that derivative, so $D[f_a(sz)]=sf'_a(sz)$.

The coordinate dictionary in Section~\ref{subsec:blocks} specifies five types:
exponentials $e^{\lambda z}$; specified first-level primitives of finite sums
$\sum_\mu q_\mu e^{\mu z}$ with $q_\mu\in\C(z)$ and arbitrary complex
additive constants, together with the finite towers $\Ein_k(\eta z)$
($\eta\ne0$, $1\le k\le h_\eta$) of \eqref{eq:ein-tower};
Bessel blocks of nonzero scale $c$ and rational order, with the four
coordinates in \eqref{eq:bessel-zero}--\eqref{eq:bessel-nonresonant};
Kummer blocks $(F,F',\Sigma,\Sigma')$ of nonzero scale $s$ and
$a\notin\Z$, defined by \eqref{eq:kummer-block}; and $L$.
The integer-order Bessel coordinates use \eqref{eq:bessel-integer} and
their derivatives; for $\nu\in\Q\setminus\Z$ the coordinates are
$(f_{\nu,c},f'_{\nu,c},f_{-\nu,c},f'_{-\nu,c})$.
The additional $(\mathrm{B2}^E)$ restrictions belong to the arithmetic statements.
A \emph{closed} list contains every lower tower level
and each second-order block's four coordinates. It adds no (B2)
derivatives, which already lie in the exponential base.

\emph{Normalization} uses \eqref{eq:bessel-integer} and
\eqref{eq:link-order} to reduce $\nu\equiv\pm\nu_0\pmod\Z$,
$0\le\nu_0\le1/2$. The endpoint $\nu_0=1/2$ reduces to exponentials;
remaining second-order representatives have $0\le\nu_0<1/2$.
Choose one scale from $\{c,-c\}$ and retain explicit expressions for the
original coordinates. Half-order rates are $\pm c$.
For the second-order differential modules $M_i$ of these blocks,
linkage means $M_j\simeq M_i\otimes\Lambda$ over $k$ for rank-one
$\Lambda$, with a twist representative $e^{\lambda z}z^\upsilon$.
The set of linkage classes is $\mathcal C$.
Theorem~\ref{thm:linkage} decides when two blocks are linked.
A finitely generated $\Gamma\subset(\C,+)$ is \emph{saturated} for the
list when it contains the listed rates, kernel rates, every $-\eta$,
Kummer scales, half-order rates $\pm c$ and linkage twist rates.
Choose a $\Z$-basis and put
\begin{equation}\label{eq:exponential-field}
\Gamma=\bigoplus_{h=1}^d\Z\lambda_h,\quad t_h=e^{\lambda_hz},\quad
T=k(t_1,\ldots,t_d),\quad t^v=\prod t_h^{v_h}.
\end{equation}
The weight of $t^v$ is $\sum_hv_h\lambda_h$.
Put $K_{\rm coord}=k(\mathcal F)$ and $K_{\rm full}=T(\mathcal F)$;
the latter is differential for a closed list and equals $K_{\rm coord}$
when the basis exponentials are listed.

First-level coordinates are primitives with rational-exponential
derivative, including $\Ein_1(\eta z)$. For
$F'=\sum q_\mu e^{\mu z}$ define the finite-support vector
\begin{equation}\label{eq:residue-vector}
\operatorname{res}(F)=(\rho_{\mu,p}(q_\mu))_{(\mu,p)},\qquad
\rho_{\mu,p}(q)=\Res_{z=p}e^{\mu(z-p)}q(z)\,dz.
\end{equation}
Here $e_{(\mu,p)}$ is the unit vector at that pair and
$\operatorname{res}(\Phi_{\mu,p})=e^{-\mu p}e_{(\mu,p)}$. Put
\begin{equation}\label{eq:rank-data}
\mathcal V_{\rm res}=\Span_\C\{\operatorname{res}(F):F\text{ first-level}\},\quad\rho=\dim\mathcal V_{\rm res},
\quad\tau=\sum_{\eta}(h_\eta-1).
\end{equation}
Thus $\tau$ counts distinct higher pairs $(\eta,k)$, $k\ge2$.
For example, $\operatorname{res}(\Ein(\eta z))=e_{(0,0)}-e_{(-\eta,0)}$.

Let $X=(X_q)$ index every listed input slot, including repeated coordinates
such as $(\eta,k)$. For each repeated coordinate $q$, let
$\operatorname{rep}(q)$ choose one representative slot. The duplicate graph
$X_q-X_{\operatorname{rep}(q)}$ is a relation in the declared polynomial ring.
A \emph{constant coordinate} is identically $c\in\C$, with graph $X_c-c$.
Subscripts name coordinates, not values; generators may use coordinate names
as variables. \emph{Selected} coordinates mean a first-level residue basis,
all distinct higher levels and one triple per linkage class, specified in
Section~\ref{sec:main-statement}.
The dimension of a prime ideal means the Krull dimension of its quotient
ring. For \emph{Laurent saturation}, clear negative powers of the
exponential variables. Then replace the resulting ideal $J$ by
$J:m^\infty=\{f:m^kf\in J\text{ for some }k\ge0\}$, where $m$ is
the product of these variables. Elimination then gives relations among the listed coordinates.
Residue constants are additive; $\mathbf c$ denotes the connection column
of Section~\ref{sec:gamma}, and unindexed $\pi$ denotes the circle constant.

\begin{theorem}\label{thm:main}
Let $\mathcal F$ be closed and normalized as above, and let $\Gamma$
be saturated for it. All first-level primitive kernels are rational-exponential
functions with complex coefficients; higher tower coordinates are defined by
$(\mathrm{B2}^\tau)$ in \eqref{eq:ein-tower}. All Bessel orders are rational.
Bessel and Kummer scales are nonzero; Kummer first parameters are nonintegral.
Put $r=\rho+\tau+3\#\mathcal C$. Then
\begin{equation}\label{eq:main-degree}
\trdeg_TK_{\rm full}=r,\qquad\trdeg_{\C(z)}K_{\rm full}=d+r.
\end{equation}
The second degree is that of $K_{\rm coord}$ when the basis exponentials
are listed. The full prime relation ideal over $T$ has dimension
$\rho+\tau+3\#\mathcal C$. Its generators include affine
first-level residue relations and one Wronskian determinant per linkage class.
They also include linear coordinate identities for linked blocks and listed exponentials.
Equality graphs identify duplicate slots of every type.
A listed constant $c$ has relation $X_c-c$.
Half-integer Bessel coordinates have the exponential graphs of
Section~\ref{sec:main-statement}, with $\pm c\in\Gamma$.
These generators and a transcendence basis are given
explicitly in Section~\ref{sec:main-statement}. Over $\C(z)$ the full
coordinate ideal is obtained by Laurent saturation and elimination.
\end{theorem}

\begin{theorem}[Normalized linkage maps]\label{thm:linkage}
All Kummer blocks here are ${}_1F_1(a;1;\cdot)$.
After normalization, $0\leq\nu_0<1/2$. For second-order blocks, a rank-one
twist is possible precisely when
\begin{align}
 \operatorname{Bess}(c,\nu_0)&\sim\operatorname{Bess}(c',\nu'_0)
 &&\Longleftrightarrow c'^2=c^2,\quad\nu'_0=\nu_0,\label{eq:link-bb}\\
 \operatorname{Kum}(s,a)&\sim\operatorname{Kum}(s',a')
 &&\Longleftrightarrow
 \begin{cases}s'=s,\ a'-a\in\Z,\\s'=-s,\ a'+a\in\Z,\end{cases}\label{eq:link-kk}\\
 \operatorname{Bess}(c,\nu_0)&\sim\operatorname{Kum}(s,a)
 &&\Longleftrightarrow\nu_0=0,\ s=\pm2c,\ a\in\tfrac12+\Z.\label{eq:link-bk}
\end{align}
A nonresonant Bessel block is not linked to a resonant one. Every such
twist has a representative $e^{\lambda z}z^\upsilon$ after a rational gauge.
Its rate is $0$ for Bessel links
and contiguity, $\pm s$ for the Kummer transformation, and $\pm c$
for the mixed link. Its power satisfies $\upsilon\in\Z$ in resonant
classes, $\upsilon\equiv0$ or $\pm2\nu_0\pmod\Z$ otherwise. Here $\upsilon$
compares the original blocks; normalized representatives have
$\upsilon\equiv0\pmod\Z$.
Section~\ref{sec:linkage} gives explicit identities for regular solutions,
normalized entire parts (E-parts) and derivative coordinates.
Their compositions obey \eqref{eq:link-composition-convention} and
\eqref{eq:link-composed}. That section fixes the constants and orientation rules.
\end{theorem}
The equivalences underlying this statement are due to Gorelov
\cite[Lemmas~12 and~14, p.~7]{Gorelov2021}; our contribution is the explicit
normalized maps.

Gorelov's rescaled, exponentially twisted identity uses
${}_1F_1(\nu+\tfrac12;2\nu+1;\cdot)$ \cite[(3), p.~260]{Gorelov2019}.
Its second parameter is one only at $\nu=0$.
Appendix~\ref{app:algorithm} refers to Supplementary Section~\ref{sup:input-model}
for generators, maps and constants.

Use Section~\ref{sec:setup}'s $E$-coordinates and $(\mathrm{B2}^E)$ primitives,
with bounds in Appendix~\ref{app:arithmetic}. Let $X$ index coordinates.
Let $\rho^E$ be the primitive residue rank, including level-one exponential integrals. Put
$\mathcal E=\operatorname{acl}_{\C}\Qbar(e^\beta:\beta\in\Qbar)$,
where algebraic closure is taken within $\C$.
Section~\ref{sec:main-statement} selects coordinates; $y_{\mathfrak c}$ is
the representative's regular first coordinate.
Independent invertible indeterminates $\Theta=(\Theta_1,\ldots,\Theta_d)$
represent basis exponentials; put $R_\Theta=\Qbar[\Theta^{\pm1}]$.
Types 2--4 are Section~\ref{sec:main-statement}'s affine residue relations,
determinants and linkage identities.

\begin{theorem}\label{thm:values}
Let $\mathcal F$ be closed and normalized, with $\Gamma$ saturated.
All rates, nodes, scales and listed constants are algebraic.
All Bessel/Kummer parameters are rational.
The coordinates are $E$-coordinates, apart from integer-order Bessel blocks
reduced to order zero by Lemma~\ref{lem:integer}.
For $\alpha\in\Qbar^\times$, put $X(\alpha)=(q(\alpha))_{q\in\mathcal F}$
and specialize only
$z=\alpha$ in Theorem~\ref{thm:main}'s generators over
$\Qbar[z^{\pm1}][\Theta^{\pm1}]$, obtaining $\mathfrak P_\alpha$ in
$R_\Theta[X]$. Suppose the basis exponentials are listed.
Define the evaluation map
\[
 \operatorname{ev}_\alpha:R_\Theta[X]\longrightarrow\C,
 \qquad \Theta_h\longmapsto e^{\lambda_h\alpha},\quad
 X_q\longmapsto q(\alpha).
\]
Let $I_\alpha=\ker(\operatorname{ev}_\alpha)$. Then
\begin{equation}\label{eq:value-ideal}
I_\alpha=\mathfrak P_\alpha,\qquad
\trdeg_{\Qbar}\Qbar(X(\alpha))=d+\rho^E+\tau+3\#\mathcal C.
\end{equation}\equationalias{eq:intro-value-degree}
The selected coordinates and basis exponentials are independent.
In particular,
\begin{equation}\label{eq:value-nonzero}
y_{\mathfrak c}(\alpha)\ne0.
\end{equation}
For several algebraic evaluation points $\alpha_i\ne0$, the same
assertions apply to the combined list of functions $q_i(\alpha_i z)$
at $z=1$, after closure, normalization and saturation. Nodes, scales
and rates are rescaled, derivative factors are retained, and all
selection conditions are tested on this combined list.
\end{theorem}

\begin{corollary}\label{cor:value-base}
Under the hypotheses of Theorem~\ref{thm:values}, write $X_{\rm ne}$ for the variables of nonexponential coordinates.
For these coordinates,
\begin{equation}\label{eq:value-exponential-base}
\trdeg_{\mathcal E}\mathcal E(X_{\rm ne}(\alpha))
=\rho^E+\tau+3\#\mathcal C.
\end{equation}
The ideal in $\mathcal E[X_{\rm ne}]$ has the nonexponential generators
of $\mathfrak P_\alpha$, with $\Theta_h=e^{\lambda_h\alpha}$.
These include types 2--4 and the half-integer exponential graphs,
together with their derivative and order-shift graphs.
They also include all duplicate graphs and $X_c-c$ for listed algebraic constants.
Equivalently, extend the full ideal to $\mathcal E[X]$ under this
substitution, then contract to $\mathcal E[X_{\rm ne}]$.
The selected values are independent over $\mathcal E$.
At $\alpha=1$, choose first-level values with jointly independent residue vectors.
Add distinct higher levels $(\eta,k)$ and one selected triple per class.
Every finite subcollection is independent over $\mathcal E$.
The degree and ideal assertions apply to each finite closed sublist.
For mixed points $q_i(\alpha_i)$, apply the conditions to the combined
rescaled list $q_i(\alpha_i z)$.
Require independent residues, distinct higher pairs $(\eta\alpha_i,k)$
and one selected triple per class.
Conditions on the unrescaled lists alone do not suffice.
\end{corollary}

For nonclosed lists, close, normalize, adjoin basis exponentials, saturate,
specialize and eliminate auxiliary coordinates, retaining duplicate/constant graphs.
The enlarged list of values need not have the same transcendence degree as the original list.

Theorem~\ref{thm:measure} in Appendix~\ref{sec:measure} gives a linear independence estimate
for $1,e,\allowbreak e^{-1},\Ein(1)$ with evaluated constants.
\setcounter{theorem}{4}
Our $H^{-60}/12$ estimate is weaker than the classical one
and applies for $H\ge H_0$.
Supplementary Section~\ref{sup:measure-comparison} compares the classical exponent and
effectiveness and gives a cruder bound for $1\le H<H_0$.

Theorem~\ref{thm:connection} gives a multi-node, mixed-family bound over
$\mathcal E$ and the regular-value base $B$.
The one-point consequence over the full $\mathcal E$ is known
\cite[Ch.~7, \S1, Theorem~1, pp.~232--233]{Shidlovskii1989};
Supplementary Section~\ref{sup:classical-comparison} gives the finite-rate reduction and distinguishes
Mahler's announcement from Rivoal's proof.
We claim no novelty for the mixed formulation over $B$.
Section~\ref{sec:gamma} states the bound and cancellation limitations.

\subsection{Relation to earlier work}
For primitive and homogeneous antecedents see Supplementary Section~\ref{sup:classical-comparison} and Appendix~\ref{app:gorelov}
\cite{Ostrowski1946,Kolchin1973,BBH1988,Gorelov2021}.
Mahler also proved that $\pi Y_0(2)/(2J_0(2))-\gamma$ is transcendental
\cite[Theorem~3 and its consequence, pp.~169--170]{Mahler1968},
a precedent for the Bessel connection row.
Fischler--Rivoal give a general relation-ideal algorithm
\cite[Theorem~4(i)--(ii), pp.~2--3 of the arXiv version]{FischlerRivoal2019}.
Supplementary Section~\ref{sup:classical-comparison} compares the regular families,
their normalizations and the restricted inhomogeneous results.

Twisted residues, integration and $\SL_2$ factors determine the pointwise
stabilizer of the coordinate field (Lemma~\ref{lem:stabilizer}).
Integral quotients and equal dimensions then identify the full ideals.
